\section{Introduction}
% What is the problem? / The Big Picture
Data has become an integral and pervasive component of modern society \cite{dykes2020effective, Zhang2020dataquilt}, influencing sectors ranging from personal decision-making to large-scale governmental policies \cite{martinezmaldonado2020layeredstorytelling, Young-Ho2021, Loerakker2024, hwang2021evidence}. Tasks involving the incorporation of data insights can vary in complexity, from laypersons using data-driven information to learn about new topics \cite{Ferreira2021climate, Zhang2020dataquilt} or about themselves \cite{martinezmaldonado2020layeredstorytelling, Young-Ho2021, Loerakker2024, yen2020,mejia2023dashboards}, to higher-order analysis by data scientists generating insights and visualisations \cite{yang2022, lowe2020requirements}, and governments evaluating data to create policies \cite{hwang2021evidence}. However, the ability to interpret and extract data insights, particularly in tasks of varying cognitive complexity, remains a challenge for many \cite{shao2024datastorytelling, nobre2024, yang2022, ge2024}. As both the availability and complexity of data continue to increase, it is important to support individuals -- particularly those who are not data savvy -- in navigating and interpreting the data \cite{marriott2018immersive, ndukwe2020teaching}.

% What have others done to solve the problem?
One of the key roles of data visualisation is to simplify the presentation of data into clear, actionable visuals, enabling people to perform tasks of varying complexity \cite{tufte2001visual}.
Data visualisations are often intended as a communication tool to enable \textit{‘visual discovery’} of insights \cite{ryan2016visual, Wang2020cheat}. Some visualisations are designed and used to enable data \textit{exploration}, where individuals are expected to independently interpret the information and extract insights for analysis \cite{yau2013datapoints, knaflic2015storytelling}. However, this approach relies either on the patterns and insights extracted from the data being readily apparent \cite{yau2013datapoints}, or on individuals possessing the necessary data analysis skills to interpret the information \cite{echeverria2018exploratory}. These are therefore commonly referred to as \textit{exploratory} visualisations, as they are primarily designed for users with data analysis expertise who are exploring unfamiliar datasets to uncover insights \cite{Kosara2013, bravo2020immersive, martinezmaldonado2020layeredstorytelling, iliinsky2011designing}.
In contrast, \textit{explanatory} visualisations serve as an alternative approach whereby judicious design choices and visualisation techniques are applied in advance to enable non-data savvy individuals to make data-driven decisions by directly communicating insights through a data story \cite{ryan2016visual, knaflic2015storytelling}. A technique that enables this explanatory approach is \textit{data storytelling} (DS), whereby information is compressed by adding visual or textual enhancements to conventional visualisations \cite{segel2010narrative, ryan2016visual, fernandeznieto2024editor}.

% Research Gap
Theoretical researchers and visualisation experts \cite{segel2010narrative, gershon2001, Daradkeh2021, boy2015storytelling, boy2017showing} have argued that DS helps users more effectively (with greater accuracy) and efficiently (in lesser time) extract insights from data compared to conventional methods \cite{elias2013storytelling, knaflic2015storytelling, zdanovic2022influence, segel2010narrative, gershon2001, krum2013cool, Hicks2009perceptual}. Sensemaking refers to the process of interpreting complex data to uncover patterns, relationships, and insights that might otherwise remain hidden \cite{campos2021makingsense}. DS has been claimed to facilitate this process by using explanatory narratives and visual cues to guide the user \cite{martinezmaldonado2020layeredstorytelling, srinivasan2019augmenting}. Yet, despite the theoretical arguments, empirical findings offer mixed results. Some studies have shown that DS can improve long-term recall \cite{bateman2010}, while others find no significant differences in memory retention compared to conventional visualisations \cite{zdanovic2022influence}. Evidence supporting the \textit{effectiveness} of DS in comprehension tasks comes from \citet{shao2024datastorytelling}, who found that DS visualisations indeed improves insight comprehension, especially for single insights. Yet, findings related to \textit{efficiency} are similarly inconsistent. While small-scale eye-tracking studies \cite{de2014evaluating, Pozdniakov2023how} (with fewer than 30 participants) suggest that DS effectively captures attention, and \citet{echeverria2018exploratory} found that DS-enhanced visualisations reduced interpretation time for educators, a large-scale study \cite{shao2024datastorytelling} showed no improvement in speed. These inconsistencies suggest the need for further research to explore how DS impacts both effectiveness and efficiency in a wide range of tasks that might require different cognitive skills. The primary goal of DS is to draw focus towards salient data points and make key takeaways from data accessible \cite{ryan2016visual}, thus existing studies predominantly focus on lower-order tasks. However, in practice, DS is commonly applied to complex topics that require higher-order thinking, such as politics, economics, or history \cite{seyser2018scrollytelling, morth2023scrollyvis}. This leaves a gap in understanding the impact that DS-enhancements could have on interpreting complex information and completing tasks that require higher-order thinking skills. 

% Contribution Statement / Aims
In this study, we aim to examine how DS-enhanced visualisations compare to conventional data visualisations in influencing users' understanding of tasks with varying levels of complexity. To explore this, we conducted a within-subjects controlled experiment where users are presented with tasks spanning the full range of Bloom's taxonomy \cite{krathwohl2002revision}, from low-order thinking skills (\textit{Remember/Identify}, \textit{Understand}) to higher-order tasks (\textit{Apply}, \textit{Analyse}, \textit{Evaluate}, \textit{Create}). Each participant was presented with four visualisations: two line charts and two choropleth maps, both with and without DS elements. An important methodological contribution of this work is the introduction of a novel \textit{Human-in-the-Loop-AI} task generation process. This approach integrates the automation capabilities of large language models with human validation to generate question-based tasks across different levels of Bloom's taxonomy. By systematically aligning cognitive complexity with task design, this process provides a scalable framework for evaluating the cognitive impact of visualisations and can be used by future research to study visualisation efficacy across diverse contexts.
A total of 128 users participated in the study where we addressed the following research questions: To what extent do DS-enhanced visualisations improve \textit{\textbf{effectiveness} (\textbf{RQ1}) and \textit{\textbf{efficiency}} (\textbf{RQ2}) of tasks involving data point identification, understanding, applying, analysing, and evaluating insights compared to conventional visualisations? To what extent do DS-enhanced visualisations improve the \textbf{\textit{quality of a creation task}} output compared to conventional visualisations? (\textbf{RQ3}); and how are visualisations with data storytelling elements \textbf{\textit{perceived}} in comparison to conventional visualisations in terms of preferences, usefulness and utility? (\textbf{RQ4})}

%% Findings
Our findings contribute to the growing body of research on data storytelling by providing empirical evidence on users' performance in a wide range of cognitive tasks. Key insights include:
\begin{itemize}
    \item DS-enhanced visualisations can influence how effectively lower-order tasks, such as finding data points and understanding insights, are completed. However, this influence does not necessarily translate to higher-order tasks where analysis, critical evaluation, or the ability to synthesise information into a novel form are necessary.
    \item For tasks at a higher cognitive level (such as those involving analysis, evaluation, or creation), DS-enhanced visualisations can positively impact the efficiency with which tasks are completed.
    \item Complex writing tasks that were written with DS-enhanced visualisations had less sentences categorised at the \textit{evaluation} level of Bloom's taxonomy. This suggests that individuals may be less inclined to critically engage with the narrative presented by DS-enhanced visualisations.
    \item Participants praised DS-enhancements for providing necessary context and clarity to the data visualisations. However, some individuals still preferred the minimalism and simplicity of conventional visualisations. 
\end{itemize}
